\section{Experimental Setup}

\subsection{Implementation}
The system is implemented in Rust. The G2P backend for vowel resolution is
provided by the \texttt{phonemizer} library~\cite{bernard2021}, which wraps the
eSpeak-NG speech synthesizer and is invoked as a Python subprocess only when a
vowel grapheme requires phonetic disambiguation; otherwise, the system operates
entirely within the Rust process.

\subsection{Gold Standard}

\subsubsection{Construction}
The gold standard corpus comprises 2,320 English--Filipino loanword pairs
constructed by native Filipino speakers, spanning spanning Philippine domains
where English loanwords are commonly used

Annotators applied the spelling conventions prescribed in~\cite{almario2014};
where the manual did not prescribe a specific nativization, they consulted
descriptive phonological sources to determine the most plausible form under the
five-vowel system.

\subsubsection{Inter-Annotator Agreement}
To validate corpus quality, inter-annotator agreement was measured by
presenting a random subsample of 60\% of entries was presented to a second
native-Filipino annotator. Agreement is reported as the percentage of entries
on which both annotators produced identical nativized forms.

\subsection{Evaluation Metrics}\label{ssec:evaluation_metrics}
The primary evaluation metric is Character Error Rate (CER), computed as the
Levenshtein edit distance between the system output and the gold standard
reference, normalized by the total character count of the reference across the
entire corpus:
\begin{equation}
    \text{CER} = \frac{\sum_{i} \mathrm{lev}(\hat{y}_i,\, y_i)}
    {\sum_{i} |y_i|}
    \times 100\%
\end{equation}
where \(\hat{y}_i\) is the system's nativized form of the \(i\)-th word, \(y_i\) is
the gold standard reference, \(\mathrm{lev}(\cdot,\cdot)\) denotes
character-level Levenshtein distance, and \(|y_i|\) is the character length of
the reference. CER is preferred over word-level accuracy because nativization
errors are often minor, and a
word-level metric would treat a one-character discrepancy as harshly as a
completely incorrect output.

We set the acceptance threshold at CER \(< 20\%\), meaning that on average
fewer than one in five reference characters is affected by an edit.
Additionally, because (\emph{e}, \emph{i}, \emph{y}) and (\emph{o}, \emph{u})
are interchangeable in Filipino~\cite{schachter1972}, we report CER under
relaxed conditions treating these pairs as equivalent. This provides a
linguistically realistic assessment by discounting errors native speakers would
not perceive.